\documentclass[10pt,a4paper]{amsart}
\usepackage{amsmath,amssymb,amsthm,graphicx}
\usepackage{verbatim}
\newenvironment{reference}{\comment}{\endcomment}
%\newenvironment{reference}{}{}
\newenvironment{slogan}{\comment}{\endcomment}
\newenvironment{history}{\comment}{\endcomment}
% For commutative diagrams we use Xy-pic
\usepackage[all]{xy}

% We use 2cell for 2-commutative diagrams.
\xyoption{2cell}
% Theorem environments.
%
\theoremstyle{plain}
\newtheorem{theorem}[subsection]{Theorem}
\newtheorem{proposition}[subsection]{Proposition}
\newtheorem{lemma}[subsection]{Lemma}

\theoremstyle{definition}
\newtheorem{definition}[subsection]{Definition}
\newtheorem{example}[subsection]{Example}
\newtheorem{exercise}[subsection]{Exercise}
\newtheorem{situation}[subsection]{Situation}

\theoremstyle{remark}
\newtheorem{remark}[subsection]{Remark}
\newtheorem{remarks}[subsection]{Remarks}

\numberwithin{equation}{subsection}

% Macros
%
\def\lim{\mathop{\mathrm{lim}}\nolimits}
\def\colim{\mathop{\mathrm{colim}}\nolimits}
\def\Spec{\mathop{\mathrm{Spec}}}
\def\Hom{\mathop{\mathrm{Hom}}\nolimits}
\def\Ext{\mathop{\mathrm{Ext}}\nolimits}
\def\SheafHom{\mathop{\mathcal{H}\!\mathit{om}}\nolimits}
\def\SheafExt{\mathop{\mathcal{E}\!\mathit{xt}}\nolimits}
\def\Sch{\mathit{Sch}}
\def\Mor{\mathop{\mathrm{Mor}}\nolimits}
\def\Ob{\mathop{\mathrm{Ob}}\nolimits}
\def\Sh{\mathop{\mathit{Sh}}\nolimits}
\def\NL{\mathop{N\!L}\nolimits}
\def\CH{\mathop{\mathrm{CH}}\nolimits}
\def\proetale{{pro\text{-}\acute{e}tale}}
\def\etale{{\acute{e}tale}}
\def\QCoh{\mathit{QCoh}}
\def\Ker{\mathop{\mathrm{Ker}}}
\def\Im{\mathop{\mathrm{Im}}}
\def\Coker{\mathop{\mathrm{Coker}}}
\def\Coim{\mathop{\mathrm{Coim}}}

% Boxtimes
%
\DeclareMathSymbol{\boxtimes}{\mathbin}{AMSa}{"02}

%
% Macros for moduli stacks/spaces
%
\def\QCohstack{\mathcal{QC}\!\mathit{oh}}
\def\Cohstack{\mathcal{C}\!\mathit{oh}}
\def\Spacesstack{\mathcal{S}\!\mathit{paces}}
\def\Quotfunctor{\mathrm{Quot}}
\def\Hilbfunctor{\mathrm{Hilb}}
\def\Curvesstack{\mathcal{C}\!\mathit{urves}}
\def\Polarizedstack{\mathcal{P}\!\mathit{olarized}}
\def\Complexesstack{\mathcal{C}\!\mathit{omplexes}}
% \Pic is the operator that assigns to X its picard group, usage \Pic(X)
% \Picardstack_{X/B} denotes the Picard stack of X over B
% \Picardfunctor_{X/B} denotes the Picard functor of X over B
\def\Pic{\mathop{\mathrm{Pic}}\nolimits}
\def\Picardstack{\mathcal{P}\!\mathit{ic}}
\def\Picardfunctor{\mathrm{Pic}}
\def\Deformationcategory{\mathcal{D}\!\mathit{ef}}
\usepackage{cleveref}
\crefname{originallemma}{original lemma}{original lemmas}
\Crefname{originallemma}{Original lemma}{Original lemmas}
\makeatletter
\newlabel{lemma-surjection}{{lemma-surjection}{source}{}{sourcepointer-1}{}}
\newlabel{lemma-surjection@cref}{{[originallemma][0][]lemma-surjection}{[1][0][]source}{}{}{}}
\makeatother
\crefname{originalreference}{original reference}{original references}
\Crefname{originalreference}{Original reference}{Original references}
\makeatletter
\newlabel{sheaves-lemma-directed-colimits-sections}{{sheaves-lemma-directed-colimits-sections}{source}{}{sourcepointer-2}{}}
\newlabel{sheaves-lemma-directed-colimits-sections@cref}{{[originalreference][0][]sheaves-lemma-directed-colimits-sections}{[1][0][]source}{}{}{}}
\makeatother
\title{Sheaves of sets on a space with a basis of quasi-compact opens are filtered colimits of finite coequalisers of constant sheaves}
\author{The Stacks Project Authors}
\date{stacks-project:modules.tex:local-snapshot-20261012; GNU FDL 1.2 or later}
\begin{document}
\maketitle
\noindent\textit{Editorial note.} All delivered mathematical text is the author's own, copied from the pinned Stacks Project chapter source sources/stacks-modules/pinned-original-source.tex; no statement or proof step was rewritten, reordered or removed.\par\medskip
\noindent\textit{Editorial note.} Preamble provenance: the Stacks Project keeps one single shared preamble for every chapter; this card takes it directly from the staged sources/stacks-modules/preamble.tex as three verbatim spans (lines 8-12, 14-18 and 154-217), so no derived file was written and nothing was added or redefined. The wrapper class and the base packages are supplied by the assembler.\par\medskip
\noindent\textit{Editorial note.} Bibliography provenance: neither the counted statement nor its proof contains a citation (this lane rescans the whole span before writing the card), so the card carries no bibliography and no entry was reconstructed.\par\medskip
\noindent\textit{Editorial note.} Selection provenance: this card was opened under the file-based control contract (CONTROL.md), with the mechanical pre-card gate run on its statement and proof interval before the card was written.\par\medskip
\noindent\textit{Allowed standard prerequisites (editorial).}\begin{itemize}
\item The presentation of any sheaf of sets as the target of a surjection from a coproduct of sheaves on a basis supported on members of the basis with finite values. Invoked at the start of the counted proof (source lines 3251-3254) as a pointer to a lemma of this chapter. Supplies the two coproducts whose coequaliser presents the sheaf.
\item The compatibility of sections over a quasi-compact open with filtered colimits of coproducts with injective transition maps. Invoked inside the counted proof at source lines 3268-3271 as a cross-chapter pointer. Allows the finite subsets of the two index sets to be matched up so that the coequaliser is taken over finite coproducts.
\end{itemize}\medskip
\ifdefined\hypertarget\hypertarget{sourcepointer-1}{}\fi
\noindent\textit{Source reference.} Original lemma lemma-surjection: sources/stacks-modules/pinned-original-source.tex, source lines 3251--3251. Navigation pointer only: the counted text refers to the Stacks item lemma-surjection, which belongs to another chapter of the project and is not part of this card. The pointer carries the target's own label as its number because this chapter's staged snapshot does not contain the other chapter's tag table; its content is not reproduced here.\par\medskip
\ifdefined\hypertarget\hypertarget{sourcepointer-2}{}\fi
\noindent\textit{Source reference.} Original reference sheaves-lemma-directed-colimits-sections: sources/stacks-modules/pinned-original-source.tex, source lines 3271--3271. Navigation pointer only: the counted text refers to the Stacks item sheaves-lemma-directed-colimits-sections, which belongs to another chapter of the project and is not part of this card. The pointer carries the target's own label as its number because this chapter's staged snapshot does not contain the other chapter's tag table; its content is not reproduced here.\par\medskip
\begin{lemma}
\label{lemma-filtered-colimit-constructibles}
Let $X$ be a topological space. Let $\mathcal{B}$ be a basis for the
topology of $X$ and assume that each $U \in \mathcal{B}$ is quasi-compact.
Then every sheaf of sets on $X$ is a filtered colimit of sheaves of the form
\begin{equation}
\label{equation-towards-constructible-sets}
\text{Coequalizer}\left(
\xymatrix{
\coprod\nolimits_{b = 1, \ldots, m} j_{V_b!}\underline{S_b}
\ar@<1ex>[r] \ar@<-1ex>[r] &
\coprod\nolimits_{a = 1, \ldots, n} j_{U_a!}\underline{S_a}
}
\right)
\end{equation}
with $U_a$ and $V_b$ in $\mathcal{B}$ and $S_a$ and $S_b$ finite sets.
\end{lemma}

\begin{proof}
By Lemma \ref{lemma-surjection} every sheaf of sets $\mathcal{F}$
is the target of a surjection whose source $\mathcal{F}_0$ is a coproduct
of sheaves the form $j_{U!}\underline{S}$ with $U \in \mathcal{B}$
and $S$ finite. Applying this to
$\mathcal{F}_0 \times_\mathcal{F} \mathcal{F}_0$
we find that $\mathcal{F}$ is a coequalizer of a pair of maps
$$
\xymatrix{
\coprod\nolimits_{b \in B} j_{V_b!}\underline{S_b}
\ar@<1ex>[r] \ar@<-1ex>[r] &
\coprod\nolimits_{a \in A} j_{U_a!}\underline{S_a}
}
$$
for some index sets $A$, $B$ and $V_b$ and $U_a$ in $\mathcal{B}$ and
$S_a$ and $S_b$ finite. For every finite subset $B' \subset B$
there is a finite subset $A' \subset A$ such that the coproduct
over $b \in B'$ maps into the coproduct over $a \in A'$ via both maps.
Namely, we can view the right hand side as a filtered colimit with
injective transition maps. Hence taking sections over the quasi-compact
opens $V_b$, $b \in B'$ commutes with this coproduct,
see Sheaves, Lemma \ref{sheaves-lemma-directed-colimits-sections}.
Thus our sheaf is the colimit of the cokernels of these maps
between finite coproducts.
\end{proof}

\end{document}
